%===========================================================
% 03_enhancement.tex - 图像增强
%===========================================================

\section{图像增强}

%-----------------------------------------------------------
% Frame 1: 为什么要增强
%-----------------------------------------------------------
\begin{frame}{为什么要增强图像？}
    \begin{columns}
        \column{0.5\textwidth}
        \textbf{低质量图像问题：}
        \begin{itemize}
            \item \highlight{对比度低}（灰蒙蒙一片）
            \item 亮部过曝或暗部欠曝
            \item 细节不清晰
            \item 试卷拍照偏暗/偏亮
        \end{itemize}

        \column{0.5\textwidth}
        \begin{alertblock}{增强目标}
            \begin{itemize}
                \item 提升对比度
                \item 恢复细节
                \item 改善视觉效果
                \item \highlight{方便后续处理（OCR、边缘检测）}
            \end{itemize}
        \end{alertblock}
    \end{columns}

    \vspace{0.5cm}
    \begin{block}{重要概念}
        \textbf{图像增强 $\neq$ 去噪} \\
        去噪是"清理垃圾"，增强是"优化显示"
    \end{block}
\end{frame}

%-----------------------------------------------------------
% Frame 2: 图像直方图
%-----------------------------------------------------------
\begin{frame}[fragile]{图像直方图}
    \textbf{定义：} 直方图表示图像中各像素值级别的分布情况

    \begin{columns}
        \column{0.5\textwidth}
        \begin{lstlisting}[basicstyle=\ttfamily\tiny]
import cv2
import numpy as np
import matplotlib.pyplot as plt

# 读取灰度图
img = cv2.imread('exam.jpg', 0)

# 计算直方图
hist = cv2.calcHist(
    [img],          # 图像
    [0],            # 通道
    None,           # 掩码
    [256],          # 像素值范围
    [0, 256]        # 范围
)

# 绘制直方图
plt.figure(figsize=(8, 4))
plt.plot(hist, color='black', linewidth=1.5)
plt.xlabel('像素值 (0-255)')
plt.ylabel('像素数量')
plt.title('图像直方图')
plt.xlim([0, 256])
plt.grid(True, alpha=0.3)
plt.show()
        \end{lstlisting}

        \column{0.5\textwidth}
        \begin{block}{直方图解读}
            \begin{itemize}
                \item \highlight{左偏}：图像偏暗
                \item \highlight{右偏}：图像偏亮
                \item \highlight{集中}：对比度低（灰蒙蒙）
                \item \highlight{分散}：对比度高
            \end{itemize}
        \end{block}

        \vspace{0.3cm}
        \begin{alertblock}{GUI演示程序：\code{gui\_02\_enhancement.py}}
			运行 \code{python gui\_launcher.py} 选择“增强演示”，可交互调节CLAHE参数和Gamma值
        \end{alertblock}
    \end{columns}
\end{frame}

%-----------------------------------------------------------
% Frame 3: 直方图均衡化
%-----------------------------------------------------------
\begin{frame}[fragile]{直方图均衡化}
    \textbf{原理：} 重新分布像素值，使直方图更均匀，增强对比度

    \begin{columns}
        \column{0.5\textwidth}
        \begin{block}{算法步骤}
            \begin{enumerate}
                \item 统计各像素值出现次数
                \item 计算累计分布函数 (CDF)
                \item 建立映射：$s = T(r) = (L-1) \cdot CDF(r)$
                \item 应用映射到所有像素
            \end{enumerate}
        \end{block}

        \column{0.5\textwidth}
        \begin{lstlisting}[basicstyle=\ttfamily\scriptsize]
import cv2
import numpy as np
import matplotlib.pyplot as plt

# 读取灰度图
gray = cv2.imread('exam.jpg', 0)

# 全局直方图均衡化
equalized = cv2.equalizeHist(gray)

# 显示对比效果
fig, axes = plt.subplots(2, 2, figsize=(10, 8))

axes[0,0].imshow(gray, cmap='gray')
axes[0,0].set_title('原图')
axes[0,0].axis('off')

axes[0,1].plot(cv2.calcHist([gray], [0], None, [256], [0, 256]))
axes[0,1].set_title('原图直方图')

axes[1,0].imshow(equalized, cmap='gray')
axes[1,0].set_title('均衡化后')
axes[1,0].axis('off')

axes[1,1].plot(cv2.calcHist([equalized], [0], None, [256], [0, 256]))
axes[1,1].set_title('均衡化直方图')

plt.tight_layout()
plt.show()
        \end{lstlisting}
    \end{columns}
\end{frame}

%-----------------------------------------------------------
% Frame 4: CLAHE
%-----------------------------------------------------------
\begin{frame}[fragile]{CLAHE (对比度受限自适应直方图均衡化)}
    \textbf{问题：} 全局均衡化会导致噪声放大

    \begin{columns}
        \column{0.5\textwidth}
        \begin{block}{CLAHE 特点}
            \begin{itemize}
                \item \highlight{局部处理}：将图像分成小块
                \item 对每个小块独立均衡化
                \item \highlight{限制对比度}：防止噪声放大
                \item 双线性插值：消除块效应
            \end{itemize}
        \end{block}

        \column{0.5\textwidth}
        \begin{lstlisting}[basicstyle=\ttfamily\scriptsize]
import cv2
import numpy as np

# 读取灰度图
gray = cv2.imread('exam.jpg', 0)

# 创建 CLAHE 对象
clahe = cv2.createCLAHE(
    clipLimit=2.0,    # 对比度限制阈值
    tileGridSize=(8, 8)  # 分块大小
)

# 应用 CLAHE
enhanced = clahe.apply(gray)

# 对比
cv2.imshow('Original', gray)
cv2.imshow('CLAHE', enhanced)
cv2.waitKey(0)
cv2.destroyAllWindows()
        \end{lstlisting}
    \end{columns}

    \vspace{0.3cm}
    \begin{alertblock}{参数说明}
        \code{clipLimit}: 对比度限制，值越大对比度越高 \\
        \code{tileGridSize}: 分块大小，越小局部效果越明显
    \end{alertblock}
\end{frame}

%-----------------------------------------------------------
% Frame 5: Gamma 校正
%-----------------------------------------------------------
\begin{frame}[fragile]{Gamma 校正}
    \textbf{公式：} $I_{out} = 255 \cdot (I_{in} / 255)^{\gamma}$

    \begin{columns}
        \column{0.5\textwidth}
        \begin{block}{效果解读}
            \begin{itemize}
                \item $\gamma < 1$：提亮暗部（校正偏暗图像）
                \item $\gamma = 1$：无变化
                \item $\gamma > 1$：压暗亮部（校正过曝图像）
            \end{itemize}
        \end{block}

        \column{0.5\textwidth}
        \begin{lstlisting}[basicstyle=\ttfamily\scriptsize]
import cv2
import numpy as np
import matplotlib.pyplot as plt

def gamma_correction(img, gamma):
    """Gamma 校正"""
    # 归一化到 [0, 1]
    normalized = img / 255.0
    # 应用 Gamma 变换
    corrected = np.power(normalized, gamma)
    # 恢复到 [0, 255]
    return np.uint8(corrected * 255)

# 读取图像
gray = cv2.imread('exam.jpg', 0)

# 不同 Gamma 值
gammas = [0.5, 1.0, 2.0]
results = [gamma_correction(gray, g) for g in gammas]

# 显示效果
fig, axes = plt.subplots(1, 3, figsize=(12, 4))
for i, (g, res) in enumerate(zip(gammas, results)):
    axes[i].imshow(res, cmap='gray')
    axes[i].set_title(f'$\gamma$ = {g}')
    axes[i].axis('off')
plt.show()
        \end{lstlisting}
    \end{columns}
\end{frame}

%-----------------------------------------------------------
% Frame 6: 图像锐化
%-----------------------------------------------------------
\begin{frame}[fragile]{图像锐化}
    \textbf{目的：} 增强边缘，恢复细节

    \begin{columns}
        \column{0.5\textwidth}
        \begin{block}{方法一：拉普拉斯算子}
            \begin{itemize}
                \item 二阶导数
                \item 突出灰度突变区域
                \item \code{cv2.Laplacian()}
            \end{itemize}
        \end{block}

        \begin{block}{方法二：Unsharp Masking}
            \begin{enumerate}
                \item 模糊原图
                \item 计算模糊图与原图之差
                \item 将差值加回原图
            \end{enumerate}
        \end{block}

        \column{0.5\textwidth}
        \begin{lstlisting}[basicstyle=\ttfamily\scriptsize]
import cv2
import numpy as np

gray = cv2.imread('exam.jpg', 0)

# 方法1: 拉普拉斯锐化
laplace = cv2.Laplacian(gray, cv2.CV_64F)
sharp1 = cv2.convertScaleAbs(laplace)

# 方法2: Unsharp Masking (更常用)
blur = cv2.GaussianBlur(gray, (5, 5), 0)
mask = cv2.subtract(gray, blur)
sharp2 = cv2.add(gray, mask)

# OpenCV 内置方法
kernel = np.array([[-1, -1, -1],
                   [-1,  9, -1],
                   [-1, -1, -1]])
sharp3 = cv2.filter2D(gray, -1, kernel)

cv2.imshow('Sharpened', sharp3)
cv2.waitKey(0)
        \end{lstlisting}
    \end{columns}
\end{frame}

%-----------------------------------------------------------
% Frame 7: 增强方法对比
%-----------------------------------------------------------
\begin{frame}{增强方法对比表}
    \begin{table}
        \centering
        \begin{tabular}{p{2.5cm}p{3cm}p{3cm}p{2cm}}
            \toprule
            \textbf{方法} & \textbf{原理} & \textbf{适用场景} & \textbf{注意} \\
            \midrule
            \code{equalizeHist} & 全局均衡化 & 整体偏暗/偏亮 & 噪声放大 \\
            \code{CLAHE} & 局部自适应 & 光照不均 & 参数调优 \\
            Gamma校正 & 幂律变换 & 偏暗/偏亮图像 & 选择合适$\gamma$ \\
            锐化 & 边缘增强 & 细节模糊 & 适度使用 \\
            \bottomrule
        \end{tabular}
    \end{table}

    \begin{alertblock}{实战建议}
        \begin{itemize}
            \item 试卷图像：先用 \code{CLAHE} 增强对比度
            \item 偏暗图像：尝试 $\gamma = 0.6\sim0.8$
            \item 模糊图像：适当锐化
            \item \highlight{组合使用}效果更好
        \end{itemize}
    \end{alertblock}
\end{frame}

%-----------------------------------------------------------
% Frame 8: 代码实战 - CLAHE + Gamma 完整流程
%-----------------------------------------------------------
\begin{frame}[fragile]{代码实战：CLAHE + Gamma 完整流程}
    \begin{lstlisting}[basicstyle=\ttfamily\scriptsize]
import cv2
import numpy as np
import matplotlib.pyplot as plt

def enhance_image(image_path, gamma=0.7, clip_limit=2.0):
    """
    图像增强完整流程
    TODO: 实现以下步骤
    """
    # Step 1: 读取图像
    gray = cv2.imread(image_path, 0)

    # Step 2: CLAHE 增强对比度
    clahe = cv2.createCLAHE(clipLimit=clip_limit, tileGridSize=(8, 8))
    enhanced = clahe.apply(gray)

    # Step 3: Gamma 校正
    normalized = enhanced / 255.0
    corrected = np.power(normalized, gamma)
    final = np.uint8(corrected * 255)

    return final

# 测试
if __name__ == '__main__':
    result = enhance_image('exam.jpg')

    # 显示对比
    original = cv2.imread('exam.jpg', 0)
    fig, axes = plt.subplots(1, 2, figsize=(10, 5))
    axes[0].imshow(original, cmap='gray')
    axes[0].set_title('原图')
    axes[1].imshow(result, cmap='gray')
    axes[1].set_title('增强后')
    plt.show()
    \end{lstlisting}

    \begin{alertblock}{挑战任务}
        实现一个函数，自动判断图像需要哪种增强方式
    \end{alertblock}
\end{frame}
